Сети связи -- отдельная подсистема автоматизированной системы управления (АСУ) НТТС и состоит из внутренней сети и связи с внешними источниками (рис.\ \cref{fig:ntts_link}). Они играют ключевую роль и обеспечивают взаимодействие между внутренними компонентами системы и внешними источниками данных соответственно. На рис.\ \ref{fig:sensors} представлена схема взаимодействия компонентов системы автоматизированного управления НТТС, где сети связи играют ключевую роль. Рассмотрим физический и логический уровни сетей связи отдельно, чтобы подчеркнуть их особенности и важность для НТТС \cite{Martin}.

\nomenclature{АСУ}{автоматизированная система управления\nomrefpage}

\begin{figure}[ht]
	\centerfloat{
		\includegraphics[width=14cm]{drawio/ntts_link.drawio.png}
	}
	\caption{Классификация сетей связи АСУ НТТС}\label{fig:ntts_link}
\end{figure}

Физический уровень определяет способы передачи данных через физические среды, такие как радиоволны, оптические волокна или проводные соединения. На этом уровне реализуются механизмы контроля целостности передаваемых данных в конкретном фрейме (кадре). Но не реализуются механизмы адресования, гарантий доставки данных и т.\ д.

Для внутренней сети связи используются исключительно проводные соединения в виду их надежности и стабильности работы, такие как:
\begin{itemize}
	\item \textit{универсальная последовательная шина (USB\ \footnote{\  Последовательный интерфейс для подключения периферийных устройств к вычислительной технике.})}:
	стандарт для подключения периферийных устройств (камер, датчиков и т.\ д.) к вычислительной технике (см.\ раздел\ \ref{subsec:ch1/sec2/sub1});

	\item \textit{сеть контроллеров (CAN-шина\ \footnote{\ Стандарт промышленной сети, ориентированный, прежде всего, на объединение в единую сеть различных исполнительных устройств и датчиков. Режим передачи — последовательный, широковещательный, пакетный.})}:
	является стандартом для автомобильной автоматики с 1980х годов и широко апробирован в различных вариантах использования;

	\item \textit{проводной Ethernet}:
	самый распространенный и надежный протокол передачи данных. Работает на канальном и физическом уровне модели взаимодействия открытых систем (OSI) \cite{Olifers, TCP_IP, ComputerNetwork} (рис.\ \cref{fig:osi_tcp}).
\end{itemize}

\nomenclature{USB}{Universal Serial Bus, универсальная последовательная шина\nomrefpage}
\nomenclature{OSI}{Open Systems Interconnection model, сетевая модель стека сетевых протоколов OSI/ISO\nomrefpage}
\nomenclature{CAN}{Controller Area Network, сеть контроллеров\nomrefpage}

\begin{figure}[ht]
	\centerfloat{
		\includegraphics[scale=0.3]{views/osi.jpg}
	}
	\caption{Сравнение моделей протоколов передачи данных}\label{fig:osi_tcp}
\end{figure}

Так как НТТС -- это мобильные объекты, то наиболее актуальными для связи с внешними источниками (в том числе друг с другом) являются беспроводные технологии, такие как сети мобильной связи (типа
4G\ \footnote{\ \ Поколение мобильной связи с повышенными требованиями. К четвертому поколению принято относить технологии, позволяющие осуществлять передачу данных со скоростью до 100 Мбит/с подвижным (с высокой мобильностью) и до 1 Гбит/с — стационарным абонентам (с низкой мобильностью).}, 5G\ \footnote{\ \ Поколение мобильной связи, действующее на основе стандартов телекоммуникаций, следующих за существующими стандартами 4G.} и др. (рис.\ \cref{fig:cellular})), сети Wi-Fi\ \footnote{\ \ Технология беспроводной локальной сети с устройствами на основе стандартов IEEE 802.11.}, спутниковой связи (например серии Ямал (рис.\ \cref{fig:yamal}) или Starlink (рис.\ \cref{fig:starlink})) и радиосвязь УКВ.

\nomenclature{4G}{fourth generation, четвертое поколение мобильной связи\nomrefpage}
\nomenclature{5G}{fifth generation, пятое поколение мобильной связи\nomrefpage}
\nomenclature{УКВ}{ультракороткие волны диапазона радиоволн\nomrefpage}

Области применения этих технологий в контексте НТТС:
\begin{itemize}
	\item сети мобильной связи обеспечивают высокую пропускную способность и низкую задержку, что критически важно для управления НТТС в~режиме реального времени \cite{6G};
	\item спутниковая связь способна обеспечивать коммуникацию в удаленных и малонаселенных районах, где отсутствует покрытие наземных сетей мобильной связи \cite{SystemComm}. При этом для движущихся объектов, таких как НТТС, необходимо использовать специализированные спутниковые системы, обеспечивающие стабильное соединение даже в условиях высокой мобильности;
	\item Wi-Fi имеет ограниченный радиус применения \cite{Olifers}, можно использовать в помещениях, для локальной связи друг с другом. Так же обеспечивают высокую пропускную способность и низкую задержку;
	\item УКВ целесообразно рассматривать как вид резервной связи, т.\ к. этот вид связи подразумевает низкую пропускную способность, но обеспечит передачу данных при отсутствии других видов связи.
\end{itemize}

Качественное сравнение представлено в табл.\ \ref{tab:PhysicalLevel}.

\begin{figure}[ht]
	\centerfloat{
		\includegraphics[scale=0.2]{views/cell.png}
	}
	\caption{Поколения сетей мобильной связи}\label{fig:cellular}
\end{figure}

\begin{figure}[ht]
	\centerfloat{
		\includegraphics[width=14cm]{views/yamal.jpg}
	}
	\caption{Зоны покрытия спутников сети Ямал}\label{fig:yamal}
\end{figure}

\begin{figure}[ht]
	\centerfloat{
		\includegraphics[width=14cm]{views/starlink.png}
	}
	\caption{Зоны работы спутников сети Starlink}\label{fig:starlink}
\end{figure}

\begin{table}[htbp]
	\centering
	\caption{Сравнение технологий беспроводной связи}\label{tab:PhysicalLevel}
	\resizebox{\textwidth}{!}{%
		\begin{threeparttable}
			\begin{tabular}{|c|c|c|c|c|c|}
				\hline
					\textbf{Технология} & \textbf{Покрытие} & \textbf{Надежность} & \textbf{Скорость} & \textbf{Задержка} & \textbf{Применение} \\ \hline
				\makecell{4G/5G} & \makecell{Городские\\районы} & Высокая & \makecell{1-10 Гбит/с} & \makecell{1-50 мс} & \makecell{Управление\\в реальном времени} \\ \hline
				\makecell{Спутниковая} & Глобальное & Средняя & \makecell{100-500 Мбит/с} & \makecell{20-800 мс} & \makecell{Удаленные\\районы} \\ \hline
				Wi-Fi & \makecell{Локальное} & Высокая & \makecell{3.5-9.6 Гбит/с} & \makecell{1-20 мс} & \makecell{Локальная связь} \\ \hline
				УКВ & \makecell{До 50 км} & Низкая & \makecell{до 10 Мбит/с} & \makecell{100-500 мс} & \makecell{Резервная связь} \\ \hline
			\end{tabular}
		\end{threeparttable}
	}
\end{table}

Логический уровень определяет протоколы и методы передачи данных, которые обеспечивают надежность, целостность и эффективность связи. Он основан и использует физический, но при этом абстрагирован, т.\ к. данные логического уровня инкапсулируются в кадры (фреймы) физического (рис.\ \cref{fig:tcp_udp}). На логическом уровне реализуются адресование (глобальное и локальное) и~ механизмы логического соединения между прикладными процессами, выполняющимися на разных объектах. При этом логическое соединение с точки зрения приложений выглядит как канал, непосредственно соединяющий процессы.

\begin{figure}[ht]
	\centerfloat{
		\includegraphics[width=8cm]{views/tcp_udp.png}
	}
	\caption{Инкапсуляция данных}\label{fig:tcp_udp}
\end{figure}

В контексте НТТС следует рассматривать:
\begin{itemize}
	\item протоколы адресования, которые реализуют механизмы поиска путей пакетов в сети:
		\begin{enumerate}
			\item межсетевые протоколы IP\ \footnote{\ \ Маршрутизируемый протокол сетевого уровня стека TCP/IP.},
			OSPF\ \footnote{\ \ Основан на технологии отслеживания состояния канала (link-state technology) и использующий для нахождения кратчайшего пути алгоритм Дейкстры.},
			BGP\ \footnote{\ \ Протокол динамической маршрутизации, он же eBGP, от exterior. Относится к классу протоколов маршрутизации внешнего шлюза.},
			IS-IS\ \footnote{\ \ Протокол внутренних шлюзов (Interial Gateway Protocol), стандартизированный ISO и использующийся в основном в крупных сетях провайдеров услуг.}
			и др. \cite{ComputerNetwork};

			\item единого сетевого (стека) протоколов информационного обмена (ЕСПИО) \cite{ESPIO};

			\nomenclature{IP}{Internet Protocol, протокол интернета\nomrefpage}
			\nomenclature{OSPF}{Open Shortest Path First, протокол динамической маршрутизации\nomrefpage}
			\nomenclature{BGP}{Border Gateway Protocol, протокол граничного шлюза\nomrefpage}
			\nomenclature{IS-IS}{протокол маршрутизации промежуточных систем\nomrefpage}
			\nomenclature{ЕСПИО}{единый стек протоколов информационного обмена\nomrefpage}
		\end{enumerate}

	\item транспортные протоколы:
		\begin{enumerate}
			\item MQTT\ \footnote{\ \ Упрощенный сетевой протокол, работающий поверх TCP/IP, ориентированный на обмен сообщениями между устройствами по принципу <<издатель -- подписчик>>.}:
			легковесный протокол для IoT\ \footnote{\ \ Концепция сети передачи данных между физическими объектами, оснащенными встроенными средствами и технологиями для взаимодействия друг с другом или с внешней средой.} с поддержкой QoS \cite{Munister};

			\item LoRaWAN\ \footnote{\ \ Беспроводная технология передачи небольших по объему данных на дальние расстояния, разработанная для распределенных сетей телеметрии, межмашинного взаимодействия и интернета вещей.}:
			энергоэффективный протокол для дальнодействующих сетей \cite{PerryLi};

			\item TCP:
			надежная передача с контролем перегрузок \cite{TCP_IP};

			\item UDP:
			минималистичный протокол без гарантий доставки \cite{ComputerNetwork};

			\item QUIC\ \footnote{\ \ Разработан Google для компенсации недостатков протокола TCP в связи с проблема развертывания новых реализаций TCP на существующих сетях связи. Позволяет мультиплексировать несколько потоков данных между двумя компьютерами, работая по протоколу UDP.}:
			современный протокол с мультиплексированием поверх UDP \cite{Olifers};

			\item MultipathTCP\ \footnote{\ \ Конкурент QUIC. Реализован как часть TCP.}:
			расширение TCP с поддержкой множества путей \cite{ComputerNetwork}.

			\nomenclature{MQTT}{message queuing telemetry transport, протокол передачи очередей сообщений\nomrefpage}
			\nomenclature{IoT}{internet of things, интернет вещей\nomrefpage}
			\nomenclature{QoS}{quality of service, качество обслуживания\nomrefpage}
			\nomenclature{LoRaWAN}{Low-power Wide-area Network, энергоэффективная сеть дальнего радиуса действия\nomrefpage}
			\nomenclature{UDP}{User Datagram Protocol, протокол пользовательских датаграмм\nomrefpage}
			\nomenclature{TCP}{Transmission Control Protocol, протокол управления передачей\nomrefpage}
			\nomenclature{QUIC}{Quick UDP Internet Connections, быстрые UDP подключения по IP\nomrefpage}
		\end{enumerate}
\end{itemize}

НТТС являются потребителями услуг, предоставляемых уровнем протоколов адресования и используют только интерфейс взаимодействия с сетевыми устройствами. Одновременно с этим, деятельность НТТС непрерывно связано с рисками, то очень важно понимать механизмы, обеспечивающие надежную передачу данных транспортными протоколами и возможности их использования:

\begin{itemize}
	\item \textit{контрольная сумма}: используется для обнаружения битовых ошибок в переданном пакете;

	\item \textit{таймер}: отсчет интервала ожидания и указание на его истечение. Последнее означает, что с высокой степенью вероятности пакет или его квитанция потеряны при передаче. В случае, если пакет доставляется с задержкой, но не теряется (преждевременное истечение интервала ожидания), либо происходит потеря квитанции, повторная передача приводит к дублированию пакета на принимающей стороне;

	\item \textit{порядковый номер сегмента}: используется для порядковой нумерации пакетов данных передаваемых от отправителя к получателю. Разрывы в порядковых номерах полученных пакетов позволяют получателю обнаружить потерю пакета. Одинаковые порядковые номера пакетов означают, что пакеты дублируют друг друга;

	\item \textit{подтверждение получения}: генерируется принимающей стороной и указывает передающей стороне на то, что соответствующий пакет или группа пакетов успешно приняты. Обычно подтверждение содержит порядковые номера успешно принятых пакетов. В зависимости от протокола различают индивидуальные и групповые подтверждения;

	\item \textit{отрицательное подтверждение}: используется получателем для сообщения отправителю, что пакет был получен некорректно. Отрицательное подтверждение обычно включает в себя порядковый номер пакета, который не был корректно получен;

	\item \textit{конвейеризация}: ограничивает диапазон порядковых номеров, которые могут использоваться для передачи пакетов. Групповая передача и рукопожатие позволяют значительно увеличить пропускную способность протоколов по сравнению с режимом ожидания подтверждений. Размер окна может быть рассчитан на основе возможностей приема и буферизации принимающей стороны, а также уровня загрузки сети.
\end{itemize}

Качественное сравнение представлено в табл.\ \ref{tab:LogicalLevel}.

 \begin{table}[htbp]
 	\centering
 	\caption{Сравнение протоколов на транспортном уровне}\label{tab:LogicalLevel}
 	\resizebox{\textwidth}{!}{%
 		\begin{threeparttable}
 			\begin{tabular}{|l|c|c|c|c|c|c|c|c|}
 				\hline
 				Протокол & Надежность & \makecell{Эффективность\\ использования\\ канала} & \makecell{Контрольная\\ сумма} & Таймер & \makecell{Подтверждение\\ получения} & \makecell{Отрицательное\\ подтверждение} & Конвейер & Применение \\ \hline
 				MQTT & \makecell{Частичная\\ (QoS 0-2)} & Низкая & Нет & Нет & \makecell{Да\\ (QoS \geq 1)} & Нет & Нет & \makecell{IoT-датчики,\\ телеметрия} \\ \hline
 				LoRaWAN & Средняя & Низкая & Да & Да & Опция & Нет & Нет & \makecell{Удаленные\\ сенсоры} \\ \hline
 				TCP & Высокая & Средняя & Да & Да & Да & Да & Да & \makecell{Веб, передача\\ файлов} \\ \hline
 				UDP & Низкая & \textbf{Макс.} & Опц. & Нет & Нет & Нет & Нет & \makecell{Потоковые данные,\\ VoIP} \\ \hline
 				QUIC & Высокая & Высокая & Да & Да & Да & Да & Да & \makecell{Мобильные\\ приложения} \\ \hline
 				MultipathTCP & \textbf{\makecell{Очень\\ высокая}} & Высокая & Да & Да & Да & Да & Да & \makecell{Критические\\ системы} \\ \hline
 			\end{tabular}
 		\end{threeparttable}
 	}
 \end{table}